\section{Regions}\label{sec:regions}

Is the result the same in named regions, two of lowland and two of mountains?

\paragraph{The regions.} The four boxes come from the earlier project (soaring-anomalous-transport,
\texttt{soaring/analysis/regions.py}) and are copied here. They are given in degrees as (west, east, south, north) with the bounds
included:
\begin{itemize}
\item Alps $(5.4,\,10.0,\,43.8,\,46.6)$;
\item Pyrenees $(-1.9,\,3.3,\,42.0,\,43.5)$;
\item Channel Coast $(-1.8,\,2.0,\,48.3,\,51.2)$;
\item Champagne-Lorraine $(2.1,\,6.2,\,48.0,\,50.6)$
\end{itemize}
The regions are drawn in Fig.~\ref{fig:map}. As in the earlier project, a flight belongs to a region when
its takeoff, the first fix of the flight, lies in the box \emph{and} its terrain class, by the rule of Section~\ref{sec:terrain},
agrees: plains or hills for Channel Coast and Champagne-Lorraine, low or high mountains for Alps and Pyrenees. The boxes do not
overlap, so a flight is in at most one region.

\begin{table}
\centering
\small
\renewcommand{\arraystretch}{1.15}\setlength{\tabcolsep}{4pt}
\begin{adjustbox}{max width=\textwidth}
\begin{tabular}{lrrrr}
\toprule
Discipline & Alps & Pyrenees & Channel Coast & Champagne-Lorraine \\
\midrule
para & 21\,891 (3\,967) & 1\,575 (398) & 4\,414 (401) & 3\,231 (464) \\
hang & 842 (150) & 11 (6)$^\dagger$ & 20 (6)$^\dagger$ & 29 (9)$^\dagger$ \\
\bottomrule
\end{tabular}
\end{adjustbox}
\caption{Flights (pilots) per discipline and region. $^\dagger$ fewer than 30 flights: not fitted.}\label{tab:regions-counts}
\end{table}


Table~\ref{tab:regions-counts} gives the cells. A cell with fewer than \UniMinCell{} flights is shown but not fitted, and a
comparison of two cells is made only when both are fitted. This leaves one comparison of the two disciplines, in the Alps
(\UniRegionHangAlpsFlights{} hang glider flights), and, within a discipline, the comparisons between regions of the paragliders. The
hang glider cells of the Pyrenees (\UniRegionHangPyrFlights{} flights), of Channel Coast (\UniRegionHangChannelFlights{}) and of
Champagne-Lorraine (\UniRegionHangChampagneFlights{}) are too small.

\paragraph{Para against hang, in the Alps.} The fixed population at $T=$ \UniTauFixed{} s has \UniRegNFirstAlpsAfixed{} paraglider
and \UniRegNSecondAlpsAfixed{} hang glider flights; $H=$ \UniRegHFirstAlpsAfixed{} \UniRegHBandFirstAlpsAfixed{} and
\UniRegHSecondAlpsAfixed{} \UniRegHBandSecondAlpsAfixed{}, a difference of \UniRegDHAlpsAfixed{} \UniRegDHBandAlpsAfixed.
With the longest $T=$ \UniRegTAlpsAlongest{} s it is \UniRegDHAlpsAlongest{} \UniRegDHBandAlpsAlongest, and on all flights
\UniRegDHAlpsBkilo{} (to 1000 s) and \UniRegDHAlpsBlongest{} (to the longest $T$) (Fig.~\ref{fig:regions-ph}, Table~\ref{tab:regions}).
The Alps give the same answer as the two mountain classes.

\begin{figure}
\centering
\includegraphics[width=\textwidth]{regions_para_hang.pdf}
\caption{In the Alps the hang gliders have the lower $\alpha(\tau)$ in the three designs. (a) fixed population at 1000 s,
(b) fixed population at the longest $T$, (c) all flights; panels as Fig.~\ref{fig:raw-A}.}\label{fig:regions-ph}
\end{figure}

\begin{table}
\centering
\small
\renewcommand{\arraystretch}{1.15}\setlength{\tabcolsep}{4pt}
\begin{adjustbox}{max width=\textwidth}
\begin{tabular}{lrrccc}
\toprule
Population & flights (pilots) para & flights (pilots) hang & $H$ para & $H$ hang & $\Delta H$ \\
\midrule
Alps, A, $T=1000$ s & 21\,672 (3\,942) & 837 (149) & \shortstack{0.906\\[-1pt]{\scriptsize [0.905, 0.907]}} & \shortstack{0.881\\[-1pt]{\scriptsize [0.876, 0.885]}} & \shortstack{$+$0.026\\[-1pt]{\scriptsize [$+$0.021, $+$0.030]}} \\
Alps, A, longest $T$ (9350 s) & 6\,743 (2\,051) & 169 (72) & \shortstack{0.895\\[-1pt]{\scriptsize [0.894, 0.897]}} & \shortstack{0.866\\[-1pt]{\scriptsize [0.857, 0.875]}} & \shortstack{$+$0.029\\[-1pt]{\scriptsize [$+$0.020, $+$0.038]}} \\
Alps, B, to 1000 s & 21\,672 (3\,942) & 837 (149) & \shortstack{0.906\\[-1pt]{\scriptsize [0.905, 0.907]}} & \shortstack{0.880\\[-1pt]{\scriptsize [0.876, 0.885]}} & \shortstack{$+$0.026\\[-1pt]{\scriptsize [$+$0.021, $+$0.030]}} \\
Alps, B, to longest $T$ (9350 s) & 6\,743 (2\,051) & 169 (72) & \shortstack{0.887\\[-1pt]{\scriptsize [0.885, 0.888]}} & \shortstack{0.855\\[-1pt]{\scriptsize [0.847, 0.862]}} & \shortstack{$+$0.032\\[-1pt]{\scriptsize [$+$0.024, $+$0.040]}} \\
\midrule
Pyrenees &  &  &  &  & not fitted: para 1575, hang 11 \\
Channel Coast &  &  &  &  & not fitted: para 4414, hang 20 \\
Champagne-Lorraine &  &  &  &  & not fitted: para 3231, hang 29 \\
\bottomrule
\end{tabular}
\end{adjustbox}
\caption{H of paragliders and hang gliders per region, as Table~\ref{tab:raw}.}\label{tab:regions}
\end{table}


\paragraph{Region against region.} For the paragliders (Fig.~\ref{fig:regions-between}, Table~\ref{tab:between}) the two mountain
regions differ little, \UniBetweenParaDHAlpsPyrAfixed{} \UniBetweenParaDHBandAlpsPyrAfixed{} for Alps minus Pyrenees at 1000 s, and the two lowland
regions differ little too, \UniBetweenParaDHChannelChampagneAfixed{} \UniBetweenParaDHBandChannelChampagneAfixed{}. Between a mountain and a lowland region
the difference of $H$, mountain minus lowland, is \UniBetweenParaDHAlpsChannelAfixed{} (Alps and Channel Coast),
\UniBetweenParaDHAlpsChampagneAfixed{} (Alps and Champagne-Lorraine), \UniBetweenParaDHPyrChannelAfixed{} (Pyrenees and Channel Coast) and
\UniBetweenParaDHPyrChampagneAfixed{} (Pyrenees and Champagne-Lorraine). The hang gliders have too few flights outside the Alps to be compared between regions.

\begin{figure}
\centering
\includegraphics[width=\textwidth]{regions_between.pdf}
\caption{Paragliders have a lower $H$ over the Alps and the Pyrenees than over the two lowland regions, and the two regions of
each kind are close. (a) $\alpha(\tau)$ of the four regions on the fixed population at 1000 s; (b) the difference of $H$ between
two regions, first minus second, in the four designs (filled: fixed population; open: all flights; circle: to 1000 s;
square: to the longest $T$), with 90\% bands.}\label{fig:regions-between}
\end{figure}

\begin{table}
\centering
\small
\renewcommand{\arraystretch}{1.15}\setlength{\tabcolsep}{4pt}
\begin{adjustbox}{max width=\textwidth}
\begin{tabular}{llcccc}
\toprule
Discipline & Regions & A, $T=1000$ s & A, longest $T$ & B, to 1000 s & B, to longest $T$ \\
\midrule
para & Alps -- Pyrenees & \shortstack{$+$0.010\\[-1pt]{\scriptsize [$+$0.007, $+$0.013]}} & \shortstack{$+$0.000\\[-1pt]{\scriptsize [$-$0.007, $+$0.006]}} & \shortstack{$+$0.011\\[-1pt]{\scriptsize [$+$0.007, $+$0.014]}} & \shortstack{$+$0.000\\[-1pt]{\scriptsize [$-$0.006, $+$0.007]}} \\
para & Alps -- Channel Coast & \shortstack{$-$0.035\\[-1pt]{\scriptsize [$-$0.037, $-$0.033]}} & \shortstack{$-$0.062\\[-1pt]{\scriptsize [$-$0.066, $-$0.059]}} & \shortstack{$-$0.034\\[-1pt]{\scriptsize [$-$0.036, $-$0.032]}} & \shortstack{$-$0.079\\[-1pt]{\scriptsize [$-$0.084, $-$0.074]}} \\
para & Alps -- Champagne-Lorraine & \shortstack{$-$0.041\\[-1pt]{\scriptsize [$-$0.043, $-$0.039]}} & \shortstack{$-$0.061\\[-1pt]{\scriptsize [$-$0.065, $-$0.058]}} & \shortstack{$-$0.041\\[-1pt]{\scriptsize [$-$0.043, $-$0.039]}} & \shortstack{$-$0.078\\[-1pt]{\scriptsize [$-$0.082, $-$0.074]}} \\
para & Pyrenees -- Channel Coast & \shortstack{$-$0.045\\[-1pt]{\scriptsize [$-$0.048, $-$0.042]}} & \shortstack{$-$0.063\\[-1pt]{\scriptsize [$-$0.070, $-$0.055]}} & \shortstack{$-$0.045\\[-1pt]{\scriptsize [$-$0.048, $-$0.041]}} & \shortstack{$-$0.072\\[-1pt]{\scriptsize [$-$0.080, $-$0.065]}} \\
para & Pyrenees -- Champagne-Lorraine & \shortstack{$-$0.051\\[-1pt]{\scriptsize [$-$0.055, $-$0.048]}} & \shortstack{$-$0.062\\[-1pt]{\scriptsize [$-$0.069, $-$0.055]}} & \shortstack{$-$0.051\\[-1pt]{\scriptsize [$-$0.055, $-$0.048]}} & \shortstack{$-$0.074\\[-1pt]{\scriptsize [$-$0.081, $-$0.067]}} \\
para & Channel Coast -- Champagne-Lorraine & \shortstack{$-$0.006\\[-1pt]{\scriptsize [$-$0.009, $-$0.004]}} & \shortstack{$+$0.001\\[-1pt]{\scriptsize [$-$0.002, $+$0.005]}} & \shortstack{$-$0.007\\[-1pt]{\scriptsize [$-$0.009, $-$0.004]}} & \shortstack{$+$0.001\\[-1pt]{\scriptsize [$-$0.004, $+$0.006]}} \\
\bottomrule
\end{tabular}
\end{adjustbox}
\caption{Difference of $H$ between two regions, first minus second, within a discipline, with 90\% bands. Not fitted (a cell with fewer than 30 flights): all pairs of hang gliders (6).}\label{tab:between}
\end{table}


\keep{Of the four regions only the Alps have enough hang glider flights; there $\Hpara-\Hhang=$ \UniRegDHAlpsAfixed{} at 1000 s,
as in the mountain classes. For the paragliders the regions order as their terrain: mountain minus lowland $H$ is \UniBetweenParaDHAlpsChannelAfixed{} to
\UniBetweenParaDHPyrChampagneAfixed, and two regions of one kind differ by \UniBetweenParaDHAlpsPyrAfixed{} or less in absolute value.}
